\section{Introduction}
\label{sec:intro}

Security and privacy issues in multi-user operating systems (OS) and application software have drawn significant attention from both academia and industry over the past decades. These systems enable concurrent access by multiple users, facilitating efficient sharing of computing resources at both hardware and software levels. At the hardware level, research has identified a series of side-channel attacks targeting shared CPU caches~\cite{rauscher2024idleleak, wang2019papp, liu2015last}, hardware sensors (\eg accelerometers, camera and gyroscope)~\cite{slocum2023going, templeman2012placeraider, gao2022inertiear}, DRAM~\cite{pessl2016drama,van2023dramaqueen}, interconnects~\cite{tan2021invisible,zhang2024beyond}, GPUs~\cite{naghibijouybari2018rendered,wei2020leaky}, and more. On the software side, numerous studies have revealed security vulnerabilities in public OS-provided APIs, leading to OS-level side-channel attacks on systems like iOS~\cite{zhang2018level, wang2023danger}, Android~\cite{diao2016no, zhang2023s}, and Linux~\cite{chen2023sync+, jiang2024sync+}. 

This paper explores OS-level side channels within file systems. File system synchronization is a critical function in the OS, ensuring that data is consistently written from volatile memory (\eg DRAM) to stable storage (\eg SSD or disk drive). This process is essential for data integrity, as it guarantees that recent changes are reliably saved, reducing data loss in the event of a system crash or power failure. Synchronization mechanisms are also integral to many OS optimizations, which improve system performance and resilience. To boost performance, techniques such as batching, scheduling, buffering, and caching are widely implemented~\cite{arpaci2018operating,mogul1994better,seltzer1990disk,silvers2000ubc}. To enhance durability against potential failures, mechanisms such as journaling, write-ahead logging, and copy-on-write are employed~\cite{chidambaram2013optimistic,chidambaram2012consistency,ganger1994metadata,hagmann1987reimplementing}.


As storage device capacities continue to increase into 10s of terabytes or more, it is increasingly common for multiple programs to share the same disk, whether on local machines or cloud-based servers~\cite{gao2019houdini, wang2012disk, bessani2014scfs}. File system synchronization methods like \textit{syncfs} play a crucial role in efficiently and reliably managing data in such shared environments. Unlike the more granular \textit{fsync}, which flushes changes for individual files, \textit{syncfs} operates at the file system level, flushing all pending changes across multiple files in a single call. This approach simplifies synchronization and enhances performance by reducing the overhead associated with multiple \textit{fsync} calls. While these optimizations enhance performance and durability, they also create new avenues for information leakage. 

Recent work~\cite{chen2023sync+, jiang2024sync+} proposed covert channels leveraging contention in \textit{fsync} calls on shared persistent storage. In contrast, our work identifies an orthogonal leakage source by exploiting delay patterns in \textit{syncfs}.  The type of leakage exposed by \textit{syncfs} is fundamentally different from contention-based leakages in that the leakage depends on the write behavior of the victim. 
Thus, unlike \textit{fsync}-based methods, our approach does not rely on the victim application executing any \textit{fsync} calls; the victim simply writes data to the file system.  This leakage is significantly richer than contention leakage, enabling both robust covert channels and high-resolution side-channel attacks. In particular, our threat model assumes a malicious application or process without special privileges, aiming to infer sensitive information about other applications sharing the file system. When the attacker invokes \textit{syncfs}, it flushes not only 
% \cheng{maybe just dirty pages, dirty page cache doesn't sound correct}
its own dirty pages but also the pending changes from other applications. By analyzing the timing delays of \textit{syncfs}, the attacker infers the I/O patterns of these processes. For example, \textit{syncfs} delays increase when other applications perform more writes to the unflushed page caches, such as during webpage loading, video streaming, or the launch of an Android application.


In this work, we first identify two novel leakage vectors through the \textit{syncfs} system call (illustrated in §~\ref{sec:reverse_engineering}). First, we demonstrate how \textit{syncfs} can profile I/O system call footprints, revealing the impact of various I/O operations on execution delays. Second, we show that \textit{syncfs} leakages can infer file write sizes, with latency patterns varying linearly below 4 KB and exhibiting dynamic behavior above 4 KB due to kernel optimizations. 

We then build up a covert channel attack by leveraging \textit{syncfs}-based leakage vectors (illustrated in §~\ref{sec:covert_channel}). By exploiting delay patterns of \textit{syncfs}, the attacker encodes information into binary bits, with high delays representing "1" and low delays representing "0". The synchronization mechanism relies on the Time Stamp Counter (TSC) for precise timing coordination between the sender and receiver. Experiments on various Linux file systems (\eg ext4, ext2. and xfs) and the Windows NTFS file system reveal bandwidths up to 7.61 Kbps with error rates ranging from 0.01\% to 1.9\%.


We also exploit leakage vectors of \textit{syncfs} to develop three end-to-end side-channel attacks (illustrated in §~\ref{sec:side_channel}). These attacks infer sensitive user activities, such as website visits, video streams, and application launches, on Linux and Android devices. By analyzing the timing patterns of \textit{syncfs}, our attacks reveal distinct footprints for these activities. To enhance the effectiveness of our attacks, we employ the Short-Time Fourier Transform (STFT) to extract frequency-domain features from the timing data. These features are further processed using a Convolutional Neural Network (CNN) model. For website fingerprinting, we identify unique \textit{syncfs} delay patterns caused by file writes during website loading, achieving an F1 score of over 93\% in both closed- and open-world scenarios. Similarly, video fingerprinting leverages timing differences from buffering operations, demonstrating high accuracy across streaming platforms like YouTube and Bilibili. In closed-world evaluations, our approach achieves F1 scores exceeding 95\% for YouTube and 89\% for Bilibili. Even in open-world settings, our method maintains robust performance, achieving F1 scores of 92.74\% for YouTube and 87.03\% for Bilibili platforms. On Android devices, application fingerprinting exploits write patterns during app launches, achieving an average F1 score of over 93\%. 

In our final attack, we adapt the \textit{syncfs}-based attack to containerized environments, targeting both detection and covert communication. By monitoring \textit{syncfs} delays, our container detection technique identifies mount and unmount operations associated with container startup and shutdown. These operations create distinct delay spikes, enabling precise identification of container activity. Furthermore, we construct a cross-container covert channel where separate containers communicate by encoding data into \textit{syncfs} delay patterns. Despite the added noise and virtualization overhead inherent in container environments, the attack achieves a bandwidth of 0.23 Kbps with an error rate of 2.4\%. 




% \yicheng{Highlight attacks on android}


In summary, the contributions of this paper are: 
\begin{itemize}
    \item \textit{New attack vector.} We analyze the \textit{syncfs} system call on Linux and Android, identifying timing leakages that serve as a novel attack vector.
    \item \textit{Covert channel attack.} We design and implement a fast, resilient covert channel leveraging \textit{syncfs}, demonstrating its effectiveness across Linux and Windows file systems.
    \item \textit{End-to-end side channel attacks.} We present three high-accuracy side-channel attacks to infer visited websites, streamed videos, and foreground applications, achieving over 90\% performance in precision, recall, and F1 score.
    \item \textit{Cross container attacks.} We extend the \textit{syncfs} attack to containerized environments, demonstrating both container detection and a cross-container covert channel with practical feasibility.
    
\end{itemize}



